% TOP_PROBLEM: {"problemId": "problem.gottschalk-s-surjunctivity-conjecture", "canonicalTitle": "Gottschalk's surjunctivity conjecture", "releaseRank": 491, "primaryDomain": "probability_ergodic_dynamics", "primaryDomainLabel": "Probability, ergodic theory & dynamics", "displayStatus": "open", "releaseStatus": "open"}
% !TeX program = lualatex
% ATTEMPT_STATUS: unresolved
\documentclass[11pt]{article}
\usepackage[margin=25mm]{geometry}
\usepackage{fontspec}
\setmainfont{Latin Modern Roman}
\usepackage{amsmath,amssymb,mathtools}
\usepackage{unicode-math}
\setmathfont{Latin Modern Math}
\usepackage{xurl,hyperref}
\hypersetup{colorlinks=true,urlcolor=blue}
\setcounter{secnumdepth}{0}
\setlength{\parindent}{0pt}
\setlength{\parskip}{5pt}
\setlength{\emergencystretch}{3em}
\begin{document}
\section*{\texorpdfstring{Gottschalk's surjunctivity conjecture}{Gottschalk's surjunctivity conjecture}}\label{491}
\subsection{Definitions and mathematical statement}
Let $G$ be any group and $A$ a finite alphabet. Give $A^G$ the product topology with $A$ discrete, and let $G$ act by shifts $(g\cdot x)(h)=x(g^{-1}h)$. A cellular automaton is a continuous shift-equivariant map $F:A^G\to A^G$, equivalently a finite-memory rule: for some finite $D\subset G$ and $\mu:A^D\to A$,
\[
 F(x)(g)=\mu\bigl((x(gh))_{h\in D}\bigr).
\]
Surjunctivity asserts
\[
 F\text{ injective}\quad\Longrightarrow\quad F\text{ surjective}
\]
for every $G,A,F$. The local rule is uniform over all sites. Results about non-uniform automata must therefore be related back to this uniform target rather than changing its definition. Neither finite generation nor countability of $G$ is imposed in the record.
\par\smallskip\textit{Formulation and concept references:} \href{https://arxiv.org/html/2503.23435v2}{[S1]}.

\subsection{Short English statement}
On every group, must a finite-alphabet cellular automaton that never merges two configurations also reach every configuration?
\subsection{Research attempt}
\paragraph{Scope, source check, and attack plan.}
This attempt by GPT-6 Astra Ultra was prepared on 27 September 2026. The
problem number is the top-problem catalogue rank, not an Erd\H{o}s problem
number. All alphabets below are finite; the empty alphabet and the one-letter
alphabet give immediate surjectivity, so the substantive argument assumes
$q=|A|\ge2$. The group need not initially be countable.

Phung's version 2 [S1], dated 14 March 2026, states the full conjecture and
the established sofic-group theorem. Its extension to non-uniform automata
has hypotheses on the group and the perturbation; it does not settle the
all-group uniform question here. The broader sofic result is also identified
in [E1]. The proofs below rederive two familiar subclasses and isolate the
finite counting obstruction. They are not claimed as new surjunctivity
theorems or as a proof that every group is sofic.

There are three concrete routes. Finite quotients can make the pigeonhole
principle applicable to periodic configurations. On amenable groups, a
forbidden image pattern creates a volume-order deficit that dominates the
boundary cost of a local inverse. For a general group, compactness still
produces that inverse and a finite forbidden pattern, but neither of the two
finite models is automatic. The work below makes this distinction precise
and tests whether ordinary cardinality or a finite-window count can bypass it.

\paragraph{Reduction A: only a finitely generated subgroup is involved.}
Let $D$ be a finite memory set and set $H=\langle D\rangle$. Denote by $F_H$
the automaton with the same rule on $A^H$. The left cosets $gH$ partition $G$,
and the value at $gh$ depends only on the values at $ghD\subseteq gH$.
After identifying each coset with $H$, the map $F$ is therefore the product
of copies of $F_H$:
\begin{equation}\label{a491:cosets}
 A^G\cong\prod_{gH\in G/H}A^H,
 \qquad F\cong\prod_{gH\in G/H}F_H.
\end{equation}
If $F_H$ is not injective, use two colliding configurations on one coset
and equal configurations on the others to disprove injectivity of $F$.
If $F_H$ is injective, equality of outputs implies equality on every coset.
Surjectivity also passes in both directions: restrict a preimage to one
coset for one direction, and choose a preimage on each coset for the other.
These are ordinary set-theoretic products; the latter choice is allowed
in the usual foundations. Consequently a counterexample on any group would
already give one on a finitely generated, hence countable, group.
This reduction removes a cardinality issue, not the group-theoretic difficulty.

\paragraph{Lemma 491.1: injectivity has a finite detection window.}
If $F$ is injective, there is a finite $K\subset G$, which we may enlarge
to contain the identity $e$, such that
\begin{equation}\label{a491:detection}
 F(x)|_K=F(y)|_K\quad\Longrightarrow\quad x(e)=y(e).
\end{equation}
\emph{Proof.} In the compact space $A^G\times A^G$, let
\[
 C=\{(x,y):x(e)\ne y(e)\},\qquad
 C_K=C\cap\{(x,y):F(x)|_K=F(y)|_K\}.
\]
These sets are closed. If every finite $C_K$ were nonempty, they would
have the finite intersection property, since intersections correspond to
unions of the finite sets $K$. Compactness would give a pair in their full
intersection, with $F(x)=F(y)$ and $x(e)\ne y(e)$, contradicting injectivity.
Thus one $C_K$ is empty. Equivariance then determines $x(g)$ from
$F(x)|_{gK}$ for every $g$.\hfill$\square$

Put $Y=F(A^G)$ and write $N_Y(W)$ for the number of its patterns on a finite set $W$.
For every finite $E\subset G$, choose one full input extending each of its
$q^{|E|}$ patterns. Distinct input patterns give distinct output patterns
on $EK$, by \eqref{a491:detection}. Hence
\begin{equation}\label{a491:lower}
                         N_Y(EK)\ge q^{|E|},
\end{equation}
Here $N_Y(EK)$ counts only patterns occurring in the actual image.
The inverse rule is only asserted on patterns occurring in $Y$.
Assigning arbitrary values to other patterns extends it to a local rule,
but does not prove that those other patterns have preimages under $F$.

\paragraph{Lemma 491.2: failure of surjectivity has a finite witness.}
The set $Y$ is compact, hence closed in the Hausdorff product $A^G$.
If $z\notin Y$, a basic open cylinder about $z$ misses $Y$. Thus there
are a finite nonempty $D_0\subset G$ and a pattern $p\in A^{D_0}$ such
that no $y\in Y$ has $y|_{D_0}=p$. Enlarge $D_0$ to contain $e$ if
necessary, extending $p$ arbitrarily. Equivariance forbids the translated
pattern at every $gD_0$. Write $d=|D_0|$ and $b=|D_0D_0^{-1}|$.

Suppose a finite window $W$ contains $r$ pairwise disjoint translates
of $D_0$. On each such block, at most $q^d-1$ patterns are allowed.
Even if every remaining site is free, the number of patterns is at most
\begin{equation}\label{a491:upper}
 N_Y(W)\le q^{|W|-rd}(q^d-1)^r
          =q^{|W|}(1-q^{-d})^r.
\end{equation}
Constraints involving two blocks can only reduce this count. No independence
of the actual image coordinates has been assumed.

\paragraph{Attempt B: complete the counting argument under the F\o lner condition.}
For the next proposition assume that for every finite $B\subset G$ and
every $\epsilon>0$ there exists a nonempty finite $E\subset G$ with
\[
             |Es\mathbin{\triangle}E|\le\epsilon|E|
                         \qquad(s\in B).
\]
This is the right F\o lner condition. It holds for amenable groups; it
is explicitly the extra hypothesis being used.

\textbf{Proposition 491.3.} Under this condition every injective finite-alphabet
cellular automaton is surjective.

\emph{Proof.} Suppose otherwise, and obtain $K,D_0,p$ as above. Define the
set of available centres
\[
 T=\{g\in E:gD_0\subset E\}
    =E\cap\bigcap_{d_0\in D_0}Ed_0^{-1}.
\]
Choose $E$ satisfying the F\o lner condition for all elements of
$K\cup D_0^{-1}$. A union bound gives
\[
 |EK|\le(1+|K|\epsilon)|E|,
 \qquad |T|\ge(1-d\epsilon)|E|.
\]
Greedily select disjoint sets $gD_0$ with centres $g\in T$. If $gD_0$
meets $hD_0$, then $h\in gD_0D_0^{-1}$, so each selected block excludes
at most $b$ possible centres. We obtain at least $r\ge |T|/b$ disjoint
blocks. Because $e\in K$, all of them lie inside $E\subset EK$.
Combining \eqref{a491:lower} and \eqref{a491:upper} gives
\[
 0\le (|EK|-|E|)\log q+r\log(1-q^{-d}).
\]
Let $c=-\log(1-q^{-d})>0$. The right-hand side is at most
\begin{equation}\label{a491:contradiction}
 |E|\left(|K|\epsilon\log q
                -\frac{(1-d\epsilon)c}{b}\right).
\end{equation}
Choose $\epsilon$ so small that $d\epsilon<1/2$ and
$|K|\epsilon\log q<c/(2b)$. Then this expression is strictly negative,
a contradiction.\hfill$\square$

For example, on $\mathbb Z$ with $K\subset[-R,R]$ and
$D_0\subset[-L,L]$, choose $E=[-N,N]\cap\mathbb Z$.
The extra output sites number at most $2R$, whereas at least a constant
multiple of $N$ disjoint forbidden blocks fit in $E$ for fixed $D_0$.
Taking logarithms makes the contradiction explicit. The radius $R$ may
depend on the automaton; the limiting window $N$ is chosen only after
that radius and the forbidden pattern have been fixed.

\paragraph{Attempt C: a different complete argument using finite quotients.}
A group is residually finite if every nonidentity element survives in
some finite quotient. For a finite set $E\subset G$, intersect the kernels
of finitely many finite quotients separating the elements $h^{-1}g$ with
$g,h\in E$, $g\ne h$. The resulting finite-index normal subgroup $N$
makes the quotient map $G\to G/N$ injective on $E$.

Let $P_N$ be the configurations constant on cosets of $N$. It has
$q^{[G:N]}$ elements. It is invariant under $F$: if $x$ is fixed by
all shifts in $N$, then $F(x)$ is also fixed by them, by equivariance.
Thus an injective $F$ restricts to an injective selfmap of a finite
set $P_N$ and is surjective on that set. Every finite input pattern can
be extended to some $P_N$ because the selected quotient is injective
on its support. The union of all $P_N$ is therefore dense in $A^G$.
As $Y$ is closed and contains every $P_N$, it is all of $A^G$.

This proves the residually finite case with no growth assumption. In
particular it illustrates why the F\o lner condition is a sufficient
mechanism rather than a necessary condition for this conjecture. To apply
this proof to an arbitrary group one would need suitable finite models
with enough compatible local patterns; residual finiteness cannot be
inserted without proof.

\paragraph{Disproof searches: what the tempting shortcuts actually show.}
Infinite-set cardinality is insufficient. On $\{0,1\}^{\mathbb N}$ the map
\[
 T(x)_0=0,\qquad T(x)_{n+1}=x_n
\]
is continuous, injective, and not surjective. It is not a counterexample:
the index set is not a group and there is a distinguished boundary site.
Nor can one use a site injection of $G$ into a proper subset while retaining
equivariance. If a map $\phi:G\to G$ commutes with every left translation,
then $\phi(g)=g\phi(e)$; it is right multiplication by a group element
and is bijective.

A nonamenable group defeats the boundary argument, but that observation
does not construct a counterexample. Equations \eqref{a491:lower} and
\eqref{a491:upper} only contradict one another when the expansion
$|EK|-|E|$ is sufficiently small relative to the packed pattern deficit.
Large expansion leaves numerical room for both bounds. Proving that a
particular counting proof stops is not proving that an injective nonsurjective
automaton exists.

Finite periodic experiments have the converse limitation. For each
fixed finite quotient, an injective selfmap is automatically surjective,
irrespective of any conjecture. A proposed infinite automaton that passes
finitely many such tests need not be injective on all configurations.
The density proof above uses every finite support and residual finiteness,
not an extrapolation from a numerical sample.

\paragraph{Finite verification and exact remaining gap.}
The companion exact script enumerates all 16 binary rules with memory
$\{0,1\}$ on cyclic groups of sizes 1 through 8. It checks image sizes,
injectivity, and surjectivity directly; the counts of bijective rules are
recorded with the output, rather than being used to certify any infinite
case. It also checks the translated-block packing inequality on cyclic
groups by exhaustive finite set comparisons. These tests target the local
conventions and the finite counting steps. Compactness and the F\o lner
limit are justified by the proofs, not by the enumeration.

\textbf{UNRESOLVED.} The strongest proved result here is surjunctivity
for groups with the displayed F\o lner property and, separately, for
residually finite groups, together with the reduction to finitely generated
groups. These belong to known theory. For a general finitely generated
group, the first missing step is a replacement for the negligible-boundary
or dense-periodic-pattern mechanism that rules out the simultaneous
existence of the finite inverse window $K$ and forbidden image pattern $p$.
Three specific continuations are to build approximately equivariant finite
models and control their bad sites, to prove an alternative pattern-counting
invariant resistant to expansion, or to construct an explicit rule whose
global injectivity and missed pattern can both be certified. None is
completed here; the established sofic theorem is not extended by this work.


\subsection{Sources}
\begin{itemize}
\item[S1]Xuan Kien Phung, On Gottschalk's surjunctivity conjecture for non-uniform cellular automata, arXiv:2503.23435v2, 2026.
\url{https://arxiv.org/html/2503.23435v2}.
\textit{Formulation source.}
\item[E1]Tullio Ceccherini-Silberstein, Michel Coornaert, and Hanfeng Li,
\emph{Expansive actions with specification of sofic groups, strong topological
Markov property, and surjunctivity}, Journal of Functional Analysis (2024),
author manuscript, introduction and consequences for full shifts.
\url{https://www.math.buffalo.edu/~hfli/surjunctive-final.pdf}.
\textit{Established sofic-group context; the narrower elementary arguments
in this attempt are proved separately. Consulted 27 September 2026.}
\end{itemize}
\end{document}
